% !TEX root = ../../../main.tex
% 来源：高中物理甲种本·第三册（电磁·光学·近代物理）
% 章节：交流电 / 三相交流电
% 原始文件：3.tex，行 1630-1647
% 类型：circuit
% ---
\begin{circuitikz}
	\draw (0,0) to [american inductors, L] (0,2);
	\draw (0,0) to [american inductors, L] (-2,-1)node[left]{$C$};	
	\draw (0,0) to [american inductors, L] (2,-1)node[right]{$B$};	
	\draw[european] (5,0) to [R=1] (5,2);
	\draw [european](5,0) to [R=3] (-2+5,-1);	
	\draw[european] (5,0) to [R=2] (2+5,-1);				
	\draw(0,2)node[left]{$A$}--node [above]{端线}(5,2)	;
	\draw (2,-1)--(2,-1.5)--(2+5,-1.5)--(2+5,-1);		
	\draw (-2,-1)--(-2,-2)--(-2+5,-2)--(-2+5,-1);
	\draw (0,0)--node [above]{中性线}(5,0);
	\node at (0,.1) [left]{$O$};
	\node at (5,.1) [right]{$O'$};
	
	\node at (0,-1){发电机};
	\node at (5,-1.2){负载};
	\draw [fill=black](0,0) circle(1.5pt);	\draw [fill=black](5,0) circle(1.5pt);
\end{circuitikz}%%\caption{三相四线制}
